% ----------------------------------------------------------
% Fundamentação Teórica
% ----------------------------------------------------------
\chapter{Fundamentação Teórica} \label{cha:fundamentacao}

Neste capítulo, será abordada a relação contemporânea entre cultura e tecnologia, a democratização do acesso cultural via meios digitais, o uso de aplicativos móveis no desenvolvimento cultural, a importância da geolocalização e a acessibilidade.

\section{\emph{Cultura e Tecnologia: Uma Relação Contemporânea}} 
\label{sec:fundamentacao:Cultura e Tecnologia: Uma Relação Contemporânea}

A relação entre cultura e tecnologia, ultimamente, está ficando mais forte com o surgimento das chamadas Tecnologias Digitais de Informação e Comunicação (TDIC). Essas tecnologias vêm trazendo mudanças profundas nos âmbitos sociais e culturais em nosso dia a dia. Com isso, surgem novas formas de interação, comunicação e também de consumo cultural, transformando de um jeito significativo o comportamento das pessoas e criando meio que uma nova realidade cultural \cite{barbosa2017tdic}

A "cultura tecnológica" ou "cultura digital" surge dessa integração das tecnologias digitais à vida cotidiana, especialmente entre os chamados nativos digitais. Para eles, a vida virtual é praticamente uma extensão natural do dia a dia. Então, pode-se dizer que essas tecnologias não só mudam as práticas culturais que já existiam antes, mas também aumentam, e muito, o acesso e a participação cultural. \cite{santos2021cultura-digital-bncc}.

Nesse cenário atual, a tecnologia acaba funcionando como uma ferramenta facilitadora, ajudando na divulgação e preservação da cultura, justamente porque ela permite mais acesso e interação com diferentes conteúdos culturais. Assim, aumentam as chances das pessoas consumirem e também produzirem cultura \cite{medeiros-cultura-tecnologica-nd}\cite{soster2021tecnologias}.

De acordo com o relatório "Cultura e Tecnologias no Brasil" do NIC.br, está rolando um aumento constante no uso de dispositivos conectados para atividades culturais, tipo assistir vídeos, ouvir músicas ou acessar conteúdos educativos. Esses dados mostram claramente que as TDIC estão realmente virando parte da vida cultural cotidiana dos brasileiros, inclusive entre grupos que antes não tinham tanto acesso aos meios tradicionais de produção cultural. \cite{nic-cultura-tecnologias-no-brasil}

Além disso, algumas políticas públicas já estão conversando diretamente com essa nova realidade, como o programa Cultura Viva, por exemplo. Elas mostram uma mudança real na forma como a promoção cultural é feita hoje em dia, tentando integrar meios digitais e comunitários para fortalecer a identidade e diversidade cultural brasileira. Por isso, para entender bem a cultura atualmente, é necessário perceber que as tecnologias digitais não mudam só a maneira como a cultura circula por aí, mas também afetam diretamente como ela é vivida, recriada e reinventada, todos os dias. \cite{prolam-cultura-viva-2025}

\section{Aplicativos Móveis como Ferramentas de Desenvolvimento Cultural} \label{sec:fundamentacao:Aplicativos Móveis como Ferramentas de Desenvolvimento Cultural}

Os aplicativos móveis têm ganhado cada vez mais importância como ferramentas para o desenvolvimento cultural, já que oferecem soluções práticas para registrar, divulgar e interagir com conteúdos culturais. Esses aplicativos permitem que os usuários acessem informações bem detalhadas sobre locais culturais, eventos e patrimônios históricos, facilitando bastante a valorização e divulgação da cultura local.\cite{oliveira-cultura-tecnologias-ufse}

A facilidade de acesso por meio dos dispositivos móveis ajuda bastante a ampliar o alcance dessas informações culturais, aumentando o interesse das pessoas e incentivando uma participação mais ativa e colaborativa. Essa dinâmica melhora, não só a experiência dos turistas, mas também cria uma relação mais próxima e afetiva das pessoas com os locais visitados.\cite{soster2021tecnologias}

E, considerando que o uso dos smartphones vem crescendo no Brasil – com cerca de 62\% das casas acessando a internet só por meio desses dispositivos, os aplicativos culturais se tornam ainda mais importantes para democratizar o acesso ao patrimônio cultural, superando inclusive barreiras geográficas e sociais.\cite{mcom-2023-internet}

Além disso, essas plataformas digitais incentivam o público a participar diretamente, com funcionalidades como envio de imagens, comentários ou sugestões, contribuindo pra uma construção coletiva da memória cultural. Isso reforça bastante o senso de pertencimento e a identidade coletiva, especialmente dentro das comunidades locais ou tradicionais.\cite{prolam-cultura-viva-2025}

Portanto, usar aplicativos móveis estrategicamente pode ajudar muito no fortalecimento e na divulgação da cultura. Essas ferramentas não só documentam e divulgam o patrimônio, como também criam novas formas de experimentar e recriar a cultura no meio digital.

 
\section{Geolocalização e Acessibilidade como Elementos Facilitadores} 

A geolocalização, que usa coordenadas geográficas nos dispositivos móveis, é uma ferramenta bem poderosa para promover e valorizar a cultura. A possibilidade de oferecer informações bem precisas, baseadas na localização do usuário, melhora muito a experiência cultural, permitindo uma interação mais direta e relevante com os patrimônios culturais e históricos visitados. \cite{soster2021tecnologias}

Além disso, ao usar a geolocalização, os aplicativos culturais conseguem sugerir rotas personalizadas e pontos culturais próximos, o que torna a experiência do usuário mais interessante e incentiva o descobrimento dos patrimônios locais. Essa tecnologia também cria ambientes imersivos, em que a informação cultural fica mais interessante porque está ligada diretamente ao lugar físico. \cite{oliveira-cultura-tecnologias-ufse}

Ao mesmo tempo, a acessibilidade aparece como algo essencial para democratizar a cultura digital. Tecnologias assistivas integradas aos aplicativos, como audiodescrição, legendas, navegação facilitada e interfaces adaptativas, garantem que pessoas com diferentes tipos de necessidades especiais também possam acessar plenamente os conteúdos culturais digitais. \cite{soster2021tecnologias}

De acordo com o relatório "Cultura e Tecnologias no Brasil", a falta de acessibilidade digital ainda é um dos principais problemas para incluir pessoas com deficiência nos ambientes virtuais. Por isso, esses recursos não são apenas desejáveis, mas realmente essenciais para uma proposta cultural que seja realmente inclusiva. \cite{nic-cultura-tecnologias-no-brasil}

A combinação da geolocalização com acessibilidade gera soluções tecnológicas que levam a cultura para públicos cada vez mais diversos. Essa integração ajuda a romper barreiras geográficas, físicas e comunicacionais, transformando a experiência cultural em algo mais pessoal, significativo e acessível para todos. Desse modo, esses dois elementos são fundamentais para garantir que aplicativos culturais sejam realmente eficazes, não só facilitando o acesso, mas promovendo um engajamento mais profundo e autêntico com os bens culturais disponíveis.


\section{Casos Análogos de Aplicativos com Finalidades Culturais} 

Analisar casos parecidos é bem importante para entender como os aplicativos móveis já estão sendo usados, com sucesso, para promover e desenvolver a cultura. Exemplos como o aplicativo "Vem Para Estância" mostram claramente como essas tecnologias ajudam a gerenciar e promover informações turísticas e culturais, oferecendo recursos como geolocalização, detalhes históricos, fotos e recomendações sobre lugares turísticos específicos.\cite{oliveira-cultura-tecnologias-ufse}

Outros exemplos interessantes são aplicativos que usam tecnologias avançadas, tipo realidade aumentada e gamificação. Eles criam experiências interativas e educativas, despertando mais interesse e envolvimento dos usuários com patrimônios históricos e culturais. Essas tecnologias fazem com que os usuários não apenas consumam conteúdos de forma passiva, mas participem de forma ativa, aumentando sua conexão emocional com esses conteúdos.\cite{soster2021tecnologias}

Conforme o relatório do NIC, iniciativas desse tipo têm sido bem importantes, especialmente em lugares com pouca oferta de políticas públicas presenciais, como em comunidades periféricas ou rurais. Nesses contextos, o aplicativo cultural acaba sendo a principal forma de aproximar os cidadãos do patrimônio coletivo. Isso reforça bastante a importância de integrar tecnologia e políticas culturais, facilitando o acesso, o registro e a valorização do saber popular.\cite{nic-cultura-tecnologias-no-brasil}

Esses aplicativos servem como base para criação de novas soluções, permitindo adaptar projetos semelhantes para realidades locais, com base em experiências que deram certo. Ao analisar esses exemplos, dá para perceber claramente alguns pontos que sempre aparecem, como a simplicidade da navegação, a valorização dos conteúdos locais e da identidade regional. Assim, o estudo desses casos ajuda a identificar boas práticas e elementos importantes para criar novos aplicativos eficientes, além de mostrar o impacto positivo dessas ferramentas na valorização e preservação do patrimônio cultural.


\section*{Resumo}

A fundamentação teórica explorou como as Tecnologias Digitais de Informação e Comunicação (TDIC) moldam a cultura contemporânea, destacando a emergência da cultura digital como fenômeno social e educacional. Analisou-se o papel dos aplicativos móveis como ferramentas fundamentais no desenvolvimento cultural, destacando o uso crescente de dispositivos móveis no Brasil e sua capacidade de tornar o patrimônio acessível. A importância da geolocalização e da acessibilidade foi ressaltada como meio de inclusão cultural efetiva, demonstrando como esses recursos ampliam o engajamento com espaços culturais. Por fim, estudos de caso evidenciaram como diferentes aplicativos vêm sendo utilizados em políticas culturais públicas e iniciativas locais, servindo como base e inspiração para a proposta do aplicativo deste trabalho.
